\section{Introduction}
Retrieval-augmented generation converts a language-model question answering problem into a joint modeling and retrieval problem where the model must not only reason over evidence, but also decide which evidence should be exposed in the first place. This separation becomes especially consequential for text-rich knowledge graphs. A relevant answer may be recoverable from node descriptions, relation paths, type constraints, or a combination of all three. A natural response is to replace one-shot retrieval with an adaptive retrieval agent where the controller searches broadly, follows promising graph neighborhoods, observes intermediate evidence, and then decides what to retrieve next.

This paper studies a deceptively simple question - When an adaptive graph-retrieval agent is given explicit pre-retrieval ``thinking'' and observation-level trace feedback, does retrieval actually improve? The question matters beyond leaderboard performance. It is also a prerequisite for personalization. A future personalized retriever should use compact user state before retrieval, not only after documents are fetched. However, if profile-like signals are injected into a retrieval controller without a calibrated action policy, they may change the trajectory without improving the ranked answer list.

We introduce \textbf{\tb}, a diagnostic protocol for adaptive knowledge-graph retrieval. \tb has two roles. First, it is an intervention in which the controller receives a compact query profile before acting and receives trace-conditioned hints when the trajectory exhibits symptoms such as repeated calls, zero-result calls, or budget exhaustion. Second, it is an instrumentation layer where every trajectory is converted into process counters that can be paired with final ranking metrics. The protocol is deliberately lightweight. It does not alter graph data, training labels, gold answers, ranking metrics, or the underlying search tools. This makes the comparison strict, so if profile and trace conditioning are useful, they should improve the agent under the same retrieval interface.

The empirical result is strong but not in the direction one might expect. On \stark, \tb makes the retrieval process more inspectable but reduces retrieval quality relative to the unconditioned ARK-style controller in our open-weight setup. The degradation appears in a Qwen3-14B validation ablation, in a Qwen3-30B-A3B held-out subset, and in shared-ID paired comparisons across PRIME, MAG, and Amazon. The trace analysis explains why the result is not merely noise as repeated calls and zero-result calls are associated with substantially larger negative rank deltas, and a stricter budget reduces trajectory length without recovering accuracy.



We organize the study around three research questions. \textbf{RQ1}: does profile/trace conditioning improve standard retrieval metrics when the underlying graph tools and evaluator are fixed? \textbf{RQ2}: if ranking quality changes, can trajectory counters explain the direction of the change? \textbf{RQ3}: does imposing a stricter interaction budget repair the failure, or does the system require a different action policy? These questions make the paper a controlled analysis of adaptive retrieval behavior rather than a comparison of unrelated retriever architectures.

Our contributions are threefold. First, we provide a reproducible profile/trace intervention for adaptive graph retrieval that isolates controller behavior from dataset and metric changes. Second, we report paired ranking and trace diagnostics showing when the intervention hurts. Third, we identify the operational mechanism of failure that profile and trace state are useful for diagnosis, but harmful as unconditional prompt context unless paired with a calibrated retrieval-action policy.
